% Compile from this directory with: pdflatex project_15.tex (twice).
\documentclass[11pt,a4paper]{article}
\usepackage[T1]{fontenc}
\usepackage[utf8]{inputenc}
\usepackage[margin=22mm]{geometry}
\usepackage{lmodern}
\usepackage{amsmath,amssymb}
\usepackage{enumitem}
\usepackage{xcolor}
\usepackage{microtype}
\usepackage{hyperref}
\definecolor{epflred}{HTML}{E3000B}
\hypersetup{colorlinks=true,urlcolor=epflred,linkcolor=epflred,
  pdftitle={ENG-654 Project 15: Lifting a custom\_3R Posture-Change Loop to custom\_6R}}
\setlength{\parindent}{0pt}
\setlength{\parskip}{6pt}
\setlist[itemize]{leftmargin=*,itemsep=5pt,topsep=4pt}
\urlstyle{tt}

\begin{document}
{\small\textbf{EPFL \textbar\ ENG-654}\hfill Kinematics-Grounded Motion Planning for Robots}
\vspace{5mm}

{\color{epflred}\Large\bfseries Project 15}

{\LARGE\bfseries Lifting a custom\_3R Posture-Change Loop\\to custom\_6R\par}

\textbf{Format:} Group of two students; simulation-based project.\\
\textbf{Course foundations:} Lectures 3--7; Exercises 2--3.

\section*{Motivation}
A positioning arm may execute a nonsingular posture change while its
attached wrist cannot maintain the desired tool orientation without becoming
singular or reaching a limit. Studying the complete robot reveals the gap
between a feasible wrist-center loop and a feasible full-pose motion.

\section*{Task}
\textbf{Overall objectives.} Extend a supplied nonsingular custom\_3R arm connection to the full
custom\_6R robot and determine how wrist IK, orientation, and limits affect
the resulting trajectory.

\begin{itemize}
  \item \textbf{Match the arm and full-robot models.} Verify that the custom\_6R
positioning subchain has the custom\_3R geometry and identify the wrist center.
Document any fixed transform between the custom\_3R operational point and
that center. Validate the D--H/URDF models before transferring a trajectory.

  \item \textbf{Inspect the supplied arm loop.} Verify that the supplied joint path
connects distinct arm IKs of the same wrist-center position and remains
regular in sampled checks. Plot $\det J_{A,v}$ in $(q_2,q_3)$, overlay the
path and endpoint IKs, and map the corresponding two-/four-arm-IK workspace
regions.

  \item \textbf{Build the full-pose command.} Choose a fixed desired tool orientation
and derive tool position from the wrist-center path and the actual fixed
tool transform. Repeat with one alternative fixed orientation. Do not
confuse a wrist-center loop with a tool-tip position command.

  \item \textbf{Recover continuous wrist solutions.} For each arm sample, solve the
residual wrist rotation and continue both admissible starting wrist flips.
Track angle periodicity consistently. Monitor pose residuals and the arm
and wrist factors in $\det J=\det J_{A,v}\det J_{W,\omega}$.

  \item \textbf{Compare arm and full-pose feasibility.} Identify whether each lift
completes, reaches a wrist singularity, or violates a joint limit. Compare
endpoint full poses and joint configurations, and refine any critical
interval. Explain why one unsuccessful fixed-orientation lift is not a
proof that the full robot lacks a nonsingular connection.
\end{itemize}

\newpage
\section*{Specifics and scope}
Use one supplied arm loop, two fixed tool orientations, and both
starting wrist flips where regular. The arm loop and arm IK will be provided,
so the new work is the wrist completion and validation. Report mathematical
regularity and limit-admissible feasibility separately. No general spatial
path optimizer or search for a new cuspidal geometry is required.

Check both translation and orientation residuals
for every claimed full-pose IK. Document angle periodicity and limits,
Jacobian conventions, and any scaling used for singular values. State all
fixed coordinates in reduced plots; a two-dimensional slice does not show
the entire 6R singular set. Refine decisive transitions, and distinguish
sampled evidence from a continuous-path guarantee. Collision freedom is
not implied by these kinematic checks.

\section*{Expected outcome}
Understand the custom\_3R-plus-wrist construction and explain how
positioning-arm topology interacts with wrist orientation. Deliver four
candidate-lift records where applicable, paired arm maps, full-pose closure
checks, and a clear comparison of arm versus full-robot failure mechanisms.

Submit reproducible code, model parameters and test data, the required plots,
a concise 5--6-page report, and a short simulation demonstration. Explain
both successful cases and failures; conclusions must follow the numerical
and geometric evidence rather than an expected outcome.

\section*{Resources}
The custom\_6R consists of the custom\_3R positioning arm plus a
spherical wrist. The arm is orthogonal with non-intersecting consecutive axes.
Its URDF is \path{lectures_main/assets/models/custom_6R/custom_6R_new.urdf};
the corresponding STL meshes are in that directory. Inspect wrist-axis
concurrency and base/tool transforms using the lecture website.

The exact custom\_3R is an orthogonal 3R robot with non-intersecting
consecutive joint axes. Its URDF is
\path{lectures_main/assets/models/custom_3R/custom_3R_new.urdf}; STL meshes are
in the same directory. Use the lecture website to inspect axes and frames,
verify D--H parameters, plot IKs, and animate configurations.
The course IK methods in Lectures 3--4 and the exact-model connection in
Lecture 6 provide the starting point. A validated arm-loop seed will be
supplied; adapting an idealized-atlas loop without revalidation is insufficient.

Use the lecture website to verify D--H parameters, visualize IKs and
singularities, and simulate paths. All listed paths are relative to the
repository root. Start the website following \path{lectures_main/README.md}.

\section*{Workload and collaboration}
Budget approximately 60 hours per student (120 person-hours per pair)
for this 2-ECTS project, following EPFL's 30-hours-per-credit convention.\footnote{\href{https://www.epfl.ch/education/teaching/teaching-guide-2/getting-started/design-a-course_1/}{EPFL: Design a Course, workload guidance}.}
Deduct any lectures or exercises included in those credits. Allocate roughly
20 team-hours to model validation, 50 to the core work, 30 to experiments,
and 20 to reporting. Reuse the specified IK and simulations. Share validation
and interpretation even if modelling and implementation are divided between
students. Hardware execution is outside scope.

\end{document}
